Eet Smakkelijk ipv smakelijk

De Duitse filosoof die zei: "Erst kommt das Fressen, dann kommt die Moral" is Bertolt Brecht. Deze uitspraak is afkomstig uit zijn toneelstuk Die Dreigroschenoper (De Driestuiversopera). Het benadrukt de prioriteit van basisbehoeften boven morele of ethische overwegingen.

Neervoeding ipv opvoeding

Het was een donderdag.
Donderdag, de dag van overvloed en expansie, weerspiegelt Greta’s impulsieve honger naar genot.

Als een warme deken
scheen de zon zijn steken

over het mulle zand
van het bloeiende land.

De aren wuifden mee
met de gouden diepzee

die immense velden
vol met stengels telde.

Omsingeld door bossen,
verspreidden de trossen

korrels van roggezaad
hun schatten in bloeistaat

uit de witte bloemen.
Men hoorde het zoemen

van actieve bijen,
die de vele weien

aangrenzend de akkers,
als dociele stakkers

om stuifmeel verlieten
om dan te genieten

van al het bijenbrood
dat zoveel voeding bood.

Zo werd niet verwaarloosd
de veeleisende kroost

van de bijenvorstin,
puur uit gemeenschapszin.

Net als noeste boeren
zoet de adel voeren.

In een klassiek landhuis
pruttelt op een fornuis

dit zuurverdiende vrucht
van valse arbeidszucht.

De haan zijn ochtendroep
viel dit keer in de soep.

Aan de muur eetgerei.
De lepels zij aan zij,

met vorken aan een stok.
De messen in een blok.

Het quasi keukenrek
hield alles op zijn plek.

Onder het houten dak
was de vloer nergens vlak.

Naast de droge kruiden
hingen dierenhuiden

aan de ruwe balken.
Geholpen door valken

toonden ze het ambacht
van de edele jacht.

De roofvogel verstoord.
De edelman vermoord.

Ook koperen potten
kon men hangend spotten

rondom de open haard.
Op tafel stond een taart,

geurig en smakelijk.
Het was vermakelijk

dat een happige meid,
tot haar enorme spijt,

zich niet kon bedwingen,
omdat stond te springen

haar dwingende vraatzucht.
Verworden tot een klucht,

brak ze bij al dat zoet
klunzig het diggelgoed.

Men betrapte haar weer.
Geen pardon deze keer.

Ze werd niet meer gehuurd
en de laan uitgestuurd.

Door “rooie poepchinees”
werd de al vroege wees,

door adellijke jeugd
als gedulde ondeugd,

uitlachend uitgezwaaid.

Alsof de weergoden
haar geen gratie boden,

vergezelde een bui
en valse rijkelui

haar aftocht met modder
die als losse flodder

door de koets zijn wielen
over haar schort vielen.

Greta, heette de meid,
ze verloor tot haar spijt,

steeds weer de zware strijd,
tegen haar gulzigheid.

Ze was mollig en rood,
oogde jeugdig, maar groot.

Haar krullen dansten vrij,
als vlammen, zij aan zij.

De lach, steeds op haar lip,
verborg de sterke grip

van een vurige brand
die naar lekkers verlangt.

Het werd ververst door smart.
De modder op haar hart

bereikte haar ogen
die niet langer logen.

Uit haar hemelse blauw
stroomden daarom al gauw

tranen uit vals beklag
over Greta’s ontslag.

Ze dwaalde werkloos rond
tot ze zoetigheid vond,

druppelend in een stroom
uit een machtige boom.

Goud geschikt in een raat,
dreven haar tot een daad

als een aap die men vangt
indien een kokos hangt

gevuld met een banaan
aan een lage liaan.

Begeerte of vrijheid?
De aap raakt beiden kwijt!

Greta’s gewoontedaad
maakte de bijen kwaad.

Haar gulzige streken
brachten bijensteken

tot achterin haar keel.
Daar werd het haar te veel.

Daar verdween de muilkorf
van Statler en Waldorf.

Hun felle commentaar
staken dieper in haar

als honende stemmen
zonder sterke remmen.

Ontgaand in hun venijn
dat ze zelf muppets zijn.

Juist die harde kritiek
maakte haar mentaal ziek.

Gaf boosheid en verdriet.
Maar ook dat mocht weer niet.

En in deze spagaat
werd haar rebel zo kwaad

dat het verstand verloor.
Ze gaf willoos gehoor

aan de vertrouwde troost.
Net als Sia weer toost

op slierend zwelg-vertier
vanuit haar chandelier

als antwoord op schaamte.
Het is het geraamte

van de humane geest
die verandering vreest.

Het raamwerk als kopie
ingebouwd in AI

toont bij complexiteit
dito eigenwijsheid.

Ieder neuraal netwerk
heeft blijkbaar als kenmerk

dat het zichzelf verstrakt
en zijn veerkracht verzwakt

naarmate het aanwast
en zich daarop aanpast.

Vandaar dat een ezel,
zo bang als een wezel,

zich vaak slechts een keer stoot.
Echter, een mens in nood

vereist veel meer geduld
en geeft de steen de schuld.

Des te meer gulzigheid,
des te hoger de spijt.

De cirkel is dus rond.
Draaiend naar de afgrond

door de duizeligheid
van de gulzige meid.

De zwelling in de keel
was een gevoelig deel

dat ademen voorkwam
en bewustzijn ontnam.

Toen helderheid verdween,
werd het zwart om haar heen.

Ze greep haar houvast mis
en viel in duisternis.

U heeft nog drie weken te eten/leven

Schoon aan de haak nou ja schoon bij Greta
